% ======================================================================
\section{Implementation}
\label{sec:impl}

Everything is implemented from scratch in PyTorch (about 1{,}800 lines) and run on an Apple M5 laptop through
the MPS (Metal) backend. The code is in the private repository linked on the first page.
\begin{itemize}
\item \texttt{treenet/trunk.py}: the trunk, its supervised training (40 epochs, SGD with Nesterov momentum,
one-cycle learning rate, weight decay $5\cdot10^{-4}$, random crop and flip), the random trunk, and feature
extraction. Training a trunk takes 2--3 minutes.
\item \texttt{treenet/tree.py}: \texttt{NodeNet}; \texttt{train\_node} (Adam, learning rate $3\cdot10^{-3}$,
batches of 512 sampled with replacement, class-balanced logistic loss, error-minimising threshold after every
pass, stop at zero errors or after 20 passes without improvement); \texttt{TreeTypeNetwork} with both child
variants, depth truncation and reduced-error pruning; and the spectral class hierarchy with
\texttt{HierarchicalTreeNetwork}.
\item \texttt{scripts/00}--\texttt{06}: data preparation, trunks, trees and baselines, double descent,
receptive fields, size minimisation, and figures. Every number and figure in this report is regenerated by
these scripts. Raw results are stored as JSON in \texttt{results/}.
\end{itemize}

\paragraph{Engineering choices that mattered.}
\begin{itemize}
\item \emph{Caching the features.} Because the trunk is frozen, its feature maps are computed once (half
precision, 0.37\,GB for CIFAR-10). Training a node is then a pass of a tiny network over a cached tensor. A
root node on all 30{,}000 CIFAR-10 tree images trains in about 15\,s, and deeper nodes, which see fewer samples, faster.
\item \emph{Class balancing and threshold calibration.} Deep children routinely face a few positives among
thousands of samples. Unweighted training then simply outputs the constant label. Balanced weighting alone
fixes that but fires on far too many negatives, which wrecked the faithful variant in our first attempt. Its
children kept creating more false alarms than they fixed, until the tree hit the depth cap. Moving the threshold
to the error-minimising value, which is the paper's own objective, fixed both problems.
\item \emph{Progress guarantee.} A node whose output is constant would hand its child the identical problem.
Such a node is retrained once with a larger learning rate and otherwise becomes a leaf. In all our runs this
only happened for sets of identical inputs with conflicting labels (\cref{sec:overfit}).
\item \emph{Disjoint splits.} The trunk, the tree, model selection and the test set use disjoint images, so
every accuracy reported is honest.
\end{itemize}
